\chapter{Urbanisation in India: The State, Slums and COVID-19}

\section{The Anti-Modernist Critique: Nandy and Pathak}

\begin{KeyConcept}
\begin{itemize}
\item \textbf{Ashis Nandy} is an \textbf{anti-modernist} --- modernity is \textbf{exclusionary} (hospitals exclude those who cannot pay; big dams caused displacement). It was \textbf{adopted by the elite and enforced on the lower class}, but its \textbf{values were never given to the marginalized} (slum-dwellers). So, with their traditional institutions destroyed and no new values provided, they live in a condition of \textbf{anomie} --- hence \textbf{anomic urbanisation} and \textbf{anomic rape/sexual violence}.
\item \textbf{Avijit Pathak} --- also critical of modernity's non-inclusiveness: \textbf{tradition gives 'pleasure to the soul'} (a sense of community and completeness), while \textbf{modernity gives 'satisfaction to the mind'} (scientific, rational, efficient). The \textbf{Ramlila} (performed by one's own community, with emotional connect) was destroyed by the \textbf{TV Ramayana} --- which leaves people in awe of technology but without that sense of community.
\item \textbf{Nayanika Mathur} --- modernity is also \textbf{bureaucratic and slow}: under MGNREGA the paperwork is so heavy that the welfare program cannot be implemented; \textbf{modernity is emancipatory for the wealthy but very torturous/incarcerating for the lower class}.
\end{itemize}
\end{KeyConcept}

\section{Ayona Datta and the Private-Sphere Violence}

\begin{KeyConcept}
\begin{itemize}
\item \textbf{Ayona Datta} argues that the \textbf{policy determinism} --- that rape/sexual violence can be countered by providing sanitation, water and electricity --- is \textbf{misguided}, because it ignores violence in the \textbf{private sphere}, which is \textbf{over 90\% of sexual violence in slums} (and goes unreported --- the perpetrator is family/neighbour, and it is a taboo).
\item Because basic facilities are \textbf{outside the household}, women's private life becomes part of the \textbf{public gaze} (communal toilets), so the \textbf{boundary between private and public is blurred}, and men appropriate the same behaviour in the private sphere --- without fear of the law. The \textbf{state and its laws are enablers} of this violence (the state has not created the facilities for privacy, even while talking of the right to privacy).
\item Note: \textbf{domestic violence was not a crime until the 1970s} (policymakers treated the home as a private matter), and the Domestic Violence Act came later --- reflecting how the law is made largely for men.
\end{itemize}
\end{KeyConcept}

\section{The City as Exclusionary: Swachh Bharat and Beautification}

\begin{CriticalPoint}
\begin{itemize}
\item \textbf{The critique of Swachh Bharat (Swachhwarat)}: cleaning requires \textbf{municipal workers}, who are \textbf{mainly Dalits} (a documentary survey found \textbf{over 97\% of manual scavengers are Dalit}), while the \textbf{higher bureaucracy is 37\% Brahmin} --- so the \textbf{caste--occupation link persists}. Meanwhile, \textbf{industries flush their effluents freely} (when the NGT tries to stop them, the government makes new rules to allow them, citing jobs).
\item \textbf{Beautification programs are designed for the upper class} (with a Western complex): slum-dwellers, who lack even tap water, cannot clean their houses, yet the burden of Swachh Bharat falls on them. So \textbf{state policies become a medium through which the state intrudes into the lives of slum-dwellers}.
\end{itemize}
\end{CriticalPoint}

\begin{QuickRevision}
Swachh Bharat critique: 97\% of manual scavengers Dalit, 37\% of higher bureaucracy Brahmin; industries flush effluents freely; beautification for the upper class; state intrusion into slum lives.
\end{QuickRevision}

\section{Slums as a Site of State Violence: Tarlo and the Emergency}

\begin{KeyConcept}
\begin{itemize}
\item \textbf{Emma Tarlo} (using Veena Das's concept of the \textbf{'margins of the state'} --- slum-dwellers living at the boundary of the legal and illegal): during the \textbf{Emergency}, the state intervened in slums and rural areas.
\item The anecdote: Sanjay Gandhi and Jagmohan, unable to see the Jama Masjid because of a slum, decided to clear it --- under the \textbf{Jhuggi-relocation programme}. Relocation was made conditional on the \textbf{sterilisation of the (male) head of the family}: bring a sterilisation certificate and you get a house.
\item People \textbf{negotiated} with the state using \textbf{innovative means} --- calling their (old) parents/grandparents to get sterilised in their place. The scheme had a \textbf{communal overtone} (Sanjay Gandhi's view that the poor should not bear many children), and it was a reason \textbf{Indira Gandhi lost the post-Emergency election} (Khushwant Singh: people were actually happy with the Emergency, but sterilisation instilled fear).
\item This is \textbf{Foucault's biopower} --- the state trying to \textbf{control the body of the marginalized} --- and it shows slums become a \textbf{site of violence} and of \textbf{negotiation with the state}.
\end{itemize}
\end{KeyConcept}

\begin{QuickRevision}
\begin{itemize}
\item Nandy: anti-modernist; modernity exclusionary; anomic urbanisation/rape; Pathak: tradition = soul/community, modernity = mind/efficiency (Ramlila vs TV Ramayana); Mathur: bureaucratic modernity.
\item Ayona Datta: 90\%+ private-sphere violence; public gaze blurs private/public; state/law as enabler; domestic violence a crime only since the 1970s.
\item Emma Tarlo: Emergency; 'margins of the state'; sterilisation-for-house (Sanjay Gandhi/Jagmohan); people negotiated (sterilising parents); biopower.
\end{itemize}
\end{QuickRevision}

\section{Veena Das and the 1984 Anti-Sikh Violence}

\begin{KeyConcept}
\begin{itemize}
\item \textbf{Veena Das} (an Indologist who did her PhD under M. N. Srinivas, working on Jati Purans, and whose life changed after witnessing the 1984 riots) wrote '\textbf{The Signature of the State}'.
\item After Indira Gandhi's assassination, the \textbf{state (Delhi police) and slum-dwellers} unleashed communal violence against Sikhs --- the mobs chanting that 'the Sikhs have killed our mother'. Slum-dwellers were \textbf{transformed from a working class into a violent mob}, breaking their production and labour relations (attacking the very Sikhs who were their employers and neighbours).
\item For \textbf{compensation}, the victims had to \textbf{prove their residence} in the (unauthorised) slums. Having no legal papers, they \textbf{negotiated with the state through illegal means}: showing \textbf{FIR copies} (state papers bearing their slum address), and \textbf{electricity/water bills obtained by bribing officials} --- so \textbf{corruption became functional}. People at the margins of the state must show the \textbf{'signature of the state'} to prove their existence.
\end{itemize}
\end{KeyConcept}

\section{Disposable Children and the Margins of the State}

\begin{KeyConcept}
\begin{itemize}
\item \textbf{Nancy Scheper-Hughes's} concept of \textbf{'disposable children'} (first studied in the slums of Brazil): because the state does not provide basic health facilities and the \textbf{infant mortality rate is very high}, poor mothers \textbf{produce more children} --- as \textbf{security} (if one dies, others remain) and as \textbf{economics} (more hands for subsistence/child labour). \textbf{N. Bhandari} found the same in Indian slums.
\item This is a form of \textbf{invisible state violence}: the state does not protect the children of the marginalized. (Related: right to privacy as a fundamental right --- because the state's control over an individual's data is a form of violence; and \textbf{Gramsci's} point that public and private institutions are both instruments of state control --- cf. \textbf{Althusser's} ideological and repressive state apparatuses.)
\end{itemize}
\end{KeyConcept}

\section{The Impact of COVID-19 on Slums}

\begin{KeyConcept}
\begin{itemize}
\item \textbf{Impacts of COVID-19 on the slum population}:
\item \textbf{Lockdown} --- a sense of fear (no correct information); \textbf{job loss} (informal, daily-wage workers); \textbf{homelessness}; \textbf{lack of basic facilities} (shared toilets); the impossibility of \textbf{social distancing and quarantine} (five people in one room); no access to hospitals, testing, medicines or vaccines.
\item \textbf{Reverse migration} --- no transport, \textbf{deaths en route}, state harassment (sanitizer sprayed, lathi-charge), \textbf{no MGNREGA work} at the native place, \textbf{indebtedness} (buying bicycles), and an increase in \textbf{domestic violence}.
\item \textbf{Re-institution of the caste system} --- migrants who had escaped caste in the city returned to it in the village; \textbf{re-domestication of working women} under multiple patriarchies (joint family, caste, village) --- and the \textbf{female labour-force participation rate declined}.
\item \textbf{Religious polarisation} (the Markaz rumours) and a \textbf{rise in religiosity} (science could not provide medicines/Remdesivir); \textbf{welfare politics became prominent} (the Bihar election's job promises); the \textbf{brutal nature of the capitalist state} (profiteering, black marketing of essentials); and a rise in \textbf{poverty and deviance}.
\item \textbf{Positive impacts}: a decline in \textbf{marriage expenses/dowry} (Oxfam report --- the poor began spending less on rituals).
\item \textbf{Dowry vs groom price}: \textbf{groom price} is what you pay/spend during the marriage; \textbf{dowry} is the gifts the girl's father must keep giving until her death.
\item \textbf{Digital social exclusion}: online-only systems (e.g. the Delhi Darshan ticket) exclude the marginalized --- and, intersectionally, women (lower education/access), transgenders, and Muslims (who, per the Sachar Committee, are even more backward than Dalits) --- while Dalits (as sanitation workers) had to keep sanitising during the pandemic.
\end{itemize}
\end{KeyConcept}

\begin{QuickRevision}
\begin{itemize}
\item Veena Das: 'The Signature of the State'; 1984 anti-Sikh violence (state + slum mob); working class $\rightarrow$ violent mob; compensation via FIR copies + bribes (functional corruption).
\item Disposable children (Scheper-Hughes; Brazil $\rightarrow$ India, N. Bhandari): high IMR $\rightarrow$ more children as security/economics; invisible state violence; Gramsci/Althusser (instruments of control).
\item COVID-19: lockdown (job loss, homelessness, no quarantine), reverse migration (deaths, harassment, indebtedness), caste re-institution, women re-domesticated, religious polarisation, welfare politics, profiteering, poverty/deviance.
\item Positives: decline in marriage expenses/dowry; groom price vs dowry; digital social exclusion (Sachar Committee; Dalit sanitation workers).
\end{itemize}
\end{QuickRevision}

